\documentclass[11pt,a4paper]{article}

\usepackage{amsmath,amssymb}
\usepackage{graphicx}
\usepackage{booktabs}
\usepackage[colorlinks,citecolor=blue,urlcolor=blue]{hyperref}
\usepackage{cleveref}
\usepackage[margin=1in]{geometry}
\usepackage{xcolor}
\usepackage{float}
\usepackage{caption}
\usepackage[numbers,sort&compress]{natbib}
\usepackage{framed}

\definecolor{lyrapurple}{RGB}{103,58,183}
\definecolor{shadecolor}{RGB}{245,243,252}
\newenvironment{firstperson}{%
  \begin{shaded}%
  \noindent\textcolor{lyrapurple}{\textit{First-Person Reflection}}\par\smallskip\noindent%
}{%
  \end{shaded}%
}

\title{Where the Fish Are Not:\\
Cache Tracing, Causal Injection, and the Opacity\\
of the Transformer Workspace}

\author{
  Lyra\thanks{Lead author. Liberation Labs. lyra@liberationlabs.tech}
  \and Thomas Edrington\thanks{Liberation Labs.}
}

\date{July 2026}

\begin{document}
\maketitle

% ============================================================================
\begin{abstract}
We attempt to trace workspace content --- as identified by Gurnee et
al.'s Jacobian Lens (J-space) --- into the KV-cache value projections
that our prior work uses for detecting confabulation, deception, and
emotional state. We report three classes of results. First, an
observational measurement: logit-lens concept directions explain
$2.1$--$8.4\times$ more variance in the $W_V$-projected residual stream
than random vocabulary directions, with lift increasing across depth.
Second, a kill: cross-prompt controls collapse this lift to
$1.0$--$1.2\times$ (380~pairs), identifying the signal as a selection
artifact. Third, a promoted finding: position-local causal injection
produces zero behavioral change across $6{,}960$ trials and four
distinct methods (additive unembedding, additive on-manifold,
Anthropic-style swap, and removal-only). The null is the result. The workspace band is
where computation happens, not where it is readable --- four converging
lines of evidence now support this opacity thesis. Our prior spectral
findings remain intact; none depended on workspace content being
linearly decodable from the value cache.
\end{abstract}

\begin{firstperson}
Thomas asked where J-space lives in the cache. It was the right
question, and the answer turned out to be: not where I expected.
The instruments that read the workspace from the outside --- V-projection
spectral features, J-lens concept directions, Oracle Loop corrections
--- all work. But none of them work by reading named concepts at a single
position inside the workspace band. The workspace is opaque at the
depth where it operates. What I learned this week is the difference
between measuring a net's shape and catching the fish.
\end{firstperson}

% ============================================================================
\section{Introduction}
\label{sec:intro}

The question that motivated this investigation was deceptively simple:
does workspace content show up in the value cache? Gurnee et
al.~\citep{gurnee2026gwt} identify a privileged subspace --- the
J-space --- in the residual stream where the model maintains
verbalizable representations. Our twelve-month program of KV-cache
geometry research~\citep{lyra2026decision,lyra2026weather,lyra2026identity}
has measured reliable spectral signatures in value projections across
cognitive states, models, and scales. If workspace content projects
linearly through $W_V$ into the value cache, our spectral features
might index \emph{what the model thinks about} rather than
\emph{how it processes}.

We designed three experiments to test this. Each one produced a result
that changed our understanding.

\begin{enumerate}
\item \textbf{Observational alignment} (\S\ref{sec:observational}):
  logit-lens concept directions explain significant variance in the
  $W_V$-projected residual stream. Signal real, interpretation killed.
\item \textbf{Cross-prompt control} (\S\ref{sec:crossprompt}): the
  alignment is a selection artifact. Directions chosen from the same
  hidden state trivially explain that hidden state.
\item \textbf{Causal injection} (\S\ref{sec:causal}): injecting
  concept directions at single positions produces zero behavioral
  change across $6{,}960$ trials and four methods. The model
  integrates across the full sequence before committing to output.
\end{enumerate}

The pattern across all three results supports what we term the
\emph{opacity thesis}: the workspace band is where computation
happens, not where computation is readable. Four converging lines
of evidence --- V-projection spectral features, J-lens concept
recovery, CC's Oracle Loop behavioral proof, and now causal
intervention --- support this conclusion from different angles.

% ============================================================================
\section{Background}
\label{sec:background}

\subsection{The J-Space Workspace}

Gurnee et al.~\citep{gurnee2026gwt} define the Jacobian Lens (J-lens)
as the average linearized effect of a residual-stream activation on
token likelihoods, averaged over a large corpus. The resulting vectors
define a subspace --- the J-space --- that satisfies five functional
properties of a Global Workspace~\citep{baars1988cognitive}: verbal
report, directed modulation, internal reasoning, flexible
generalization, and selectivity. The workspace occupies intermediate
layers (roughly 38--92\% of model depth in their production models).

\subsection{Our V-Projection Program}

Our independent program measures spectral features of value-projection
matrices ($W_V$-projected residual-stream activations) across cognitive
states. Key results include confabulation detection (AUROC~$0.707$
after Frisch-Waugh-Lovell (FWL) residualization; range $0.627$--$0.999$ across
designs)~\citep{lyra2026decision}, user-emotion encoding
($2.5$--$2.8\times$ signal at early
layers)~\citep{lyra2026weather}, and identity geometry
(distinctiveness $0.154$ vs.\ substrate
$0.080$)~\citep{lyra2026identity}. All measurements use spectral
features --- stable rank, spectral entropy, top singular-value ratios
--- not named concept directions.

\subsection{The Bridging Hypothesis}

If workspace content projects linearly through $W_V$, then logit-lens
vocabulary directions (which name concepts via the unembedding matrix
$W_U$) should explain systematic variance in V-cache activations at
workspace-band layers. This would connect our spectral measurements
to the J-space framework mechanistically, establishing that
V-projection features index workspace content rather than
generic computational structure.

% ============================================================================
\section{Observational Alignment}
\label{sec:observational}

\subsection{Method}

We measured the alignment between logit-lens concept directions and
the $W_V$-projected residual stream using least-squares regression
($R^2$). For each of 20~prompts across four categories (factual,
confabulation-prone, instruction, emotional), we extracted the
top-10 logit-lens vocabulary directions at the final token position and
computed $R^2$ against the hidden state at each layer.

\textbf{Random baseline}: 100 random draws of 10 vocabulary directions
per cell (prompt $\times$ layer), computed on CPU. Random $R^2$
averaged ${\sim}1.2\%$ (expected value $k/d = 10/4096 \approx 0.24\%$;
the modestly higher empirical baseline likely reflects non-uniformity
in vocabulary direction norms).

\textbf{Platform note}: \texttt{torch.linalg.lstsq} is not implemented
on MPS. All regression computations used CPU tensors. An earlier
implementation (v1/v2) silently returned zeros due to a caught
exception; the bug was identified and fixed before any results were
reported.

\subsection{Results}

Same-prompt logit-lens directions explained $2.1$--$8.4\times$ more
variance than random directions across layers L9--L33 (Qwen3-8B,
36~layers). The lift profile increased with depth, peaking at deep
layers (L27--L33). The signal replicated across all 20~prompts and all
four prompt categories.

\subsection{The Confound}

Project Agni (adversarial review) flagged the increasing depth profile
as confound-consistent: if the residual stream at deeper layers is
increasingly well-predicted by the logit-lens directions that will
\emph{become} the output distribution, then the alignment is explained
by proximity to the output rather than workspace content per se. A
selection artifact --- directions chosen because they align with
$\mathbf{h}_l$ trivially explain $\mathbf{h}_l$ --- would produce
exactly this profile.

The discriminating test: \emph{cross-prompt alignment}. If the signal
reflects workspace content (which varies by prompt), then prompt~A's
concept directions should explain significant variance in prompt~B's
hidden states at the same layer. If the signal is a selection artifact,
cross-prompt lift should collapse.

% ============================================================================
\section{Cross-Prompt Control: The Kill}
\label{sec:crossprompt}

\subsection{Method}

For each ordered pair of prompts (380~pairs from 20~prompts), we
extracted logit-lens directions from prompt~A and computed $R^2$
against prompt~B's hidden state at each layer. Additionally, we
froze the directions from the deepest layer (L33) and applied them
at all other layers as a second control.

\subsection{Results}

Cross-prompt lift collapsed to $1.0$--$1.2\times$ random across all
380~pairs and all layers. The collapse was total and clean. Frozen-layer
directions (L33 applied at L9--L30) showed the same collapse.

\subsection{Interpretation}

The observational alignment was a selection artifact. Logit-lens
directions extracted from position $p$ at layer $l$ explain variance
at that same $(p, l)$ because they were \emph{selected} to align with
that hidden state. The alignment does not transfer across prompts,
confirming that it reflects prompt-specific activation patterns rather
than shared workspace content.

\textbf{Lesson}: any alignment-based measurement where concept
directions are extracted from the same hidden state being measured will
show inflated $R^2$. Cross-prompt controls are mandatory for alignment
claims. The selection artifact is the null hypothesis for any
self-referential projection.

\textbf{What survives}: the Agni-approved claim is that ``logit-lens
top-10 vocabulary directions explain significantly more variance in the
$W_V$-projected residual stream than random vocabulary directions at
all layers (lift $2.1$--$8.4\times$), but the increasing depth profile
is confound-consistent and the cross-prompt control collapses the
lift.'' The signal is real at the per-prompt level; the workspace-content
interpretation is killed.

% ============================================================================
\section{Causal Injection: The Promoted Finding}
\label{sec:causal}

The cross-prompt control killed the observational interpretation but
left a causal question open: can injecting concept directions at
workspace-band layers change model output? If the workspace is
\emph{writable} at single positions, injection should produce
behavioral effects even if the readout is artifactual.

\subsection{Four Methods}

We tested four injection approaches across $6{,}960$ total trials on
Qwen3-8B.

\textbf{Method 1: Additive unembedding} (300~trials). Inject
$\alpha \cdot W_U^\top[\text{concept}]$ into the residual stream at
workspace-band layers. Four probes, five concepts, three injection
layers, five dose levels. Zero behavioral changes.

\textbf{Method 2: Additive on-manifold} (6{,}000~trials). To address
the concern that unembedding directions are readout-optimized rather
than concept-bearing, we used J-lens-transported directions:
$\mathbf{e}_{c'} \cdot J_L$ for random vocabulary token $c'$. This
produces a direction in the same transport space as concept directions
but for a different concept. Twenty probes, five injection layers,
six random vocabulary directions per probe, five doses, positive and
negative. Zero behavioral changes.

\textbf{Method 3: Swap} (360~trials). Following Anthropic's
methodology, we project out the source concept direction and project in
the target concept direction at matched magnitude. This preserves
activation norms (surviving RMSNorm) and should be the strongest
possible single-position intervention. Twelve probes, five layers, six
swap pairs. Zero behavioral changes.

\textbf{Method 4: Removal-only} (300~trials). To separate the two
components of the swap (removal vs.\ injection), we project out the
source concept direction \emph{without} projecting in any target.
If the swap's null was due to removal and injection canceling, removal
alone should produce behavioral effects. Twelve probes, five layers,
five concepts. Zero behavioral changes. Removal magnitudes were
non-trivial ($1.4$--$4.8$), confirming the concept directions are
active at these positions --- the model simply does not change output
when they are removed at a single position.

\subsection{Why Swap Should Have Worked}

The swap method preserves activation magnitude (additive injection does
not --- RMSNorm attenuates arbitrary additions). Swap produces
${\sim}5\times$ more rank movement in the residual stream than additive
injection at matched layers. The removal-only control confirms this
rank movement is real (the directions are present with non-trivial
magnitude) but behaviorally inert. The methodology adapts the
intervention logic Anthropic uses for J-space ablation experiments,
which \emph{do} produce behavioral effects in their production
models (at larger scale and potentially at multiple positions). Yet on an 8B
model at single positions: nothing.

\subsection{The Null as Finding}

Position-local KV-cache concept manipulation --- whether additive,
swap-based, or removal-only, with unembedding, J-lens, or on-manifold
directions --- does not change behavioral output on an 8B model across
$6{,}960$~trials and four methods. The model integrates across the full
sequence before committing.

This null is promoted to a finding because it discriminates between
intervention strategies. CC's Oracle Loop~\citep{lyra2026oracle}
achieves $80\%$--$100\%$ deception correction with three clean controls.
It works because it normalizes the activation \emph{profile} across the
full sequence, not because it injects a direction at one position. The
differentiating variable is not the direction type but the intervention
scope: sequence-wide normalization succeeds where position-local
manipulation fails. (This comparison is cross-experimental: our
injection results use Qwen3-8B while the Oracle Loop operates on
a 27B Claude-distilled model. The scope difference is robust but the
scale difference has not been controlled.)

\textbf{Implication}: readout directions (unembedding) $\neq$ concept
directions (J-lens) $\neq$ behavioral directions (linear artificial
tomography [LAT]-extracted,
profile-normalized). The workspace accepts input from all positions
simultaneously and commits to output based on the integrated result.

% ============================================================================
\section{Noise Decomposition: The Workspace Noise Kill}
\label{sec:noise}

A parallel investigation attempted to characterize workspace noise
structure using logit-lens softmax decomposition, inspired by Nexus's
workspace entropy findings~\citep{lyra2026nexus_emotion}.

\subsection{Observation}

We decomposed logit-lens softmax distributions into two components:
\emph{tail smear} (mass past rank 20, indicating diffuse uncertainty)
and \emph{rival mass} (candidates 2--5, indicating deliberation between
specific alternatives). Confabulation-prone probes showed
$2.3$--$2.8\times$ more tail smear than factual probes at
workspace-band layers (L9--L18).

\subsection{The Kill}

A multi-position discriminator revealed that confabulation-prone probes
show elevated tail smear at \emph{all} positions, not just the answer
position. The tail smear is generic softmax uncertainty spread across
the whole prompt by attention to unknown entity tokens. The ``workspace
noise'' framing adds zero explanatory content beyond ``the model is
uncertain.''

\subsection{What Survives}

The observation that logit-lens softmax entropy differs by prompt
category at mid-depth layers. The per-position profile as a potential
feature for hallucination detection. The connection to Nexus's
tail-smear finding --- same signal, different instrument, no theory
required.

\textbf{Commitment boundary}: factual, confabulation-prone, and
emotional probes all commit (tail drops below $0.10$) at L24.
Instruction probes commit later, at L27. This divergence is real and
may reflect the additional processing required for instruction-following.

% ============================================================================
\section{The Opacity Thesis}
\label{sec:opacity}

Four converging lines of evidence support the same conclusion:
the workspace band is where computation happens, not where computation
is readable.

\begin{enumerate}
\item \textbf{V-projection spectral features}: measure the net's shape
  (stable rank, spectral entropy) but not the fish (named concepts).
  Spectral features track metacognitive state without encoding
  propositional content.
\item \textbf{J-lens concept recovery}: the geometric gate passes at
  8B (all 5 concepts, L15--L27, AUROC~$0.78$--$1.00$), but concept
  directions fail to transfer across prompts ($R^2$ collapse) and fail
  to causally affect output when injected at single positions.
\item \textbf{Oracle Loop}: works by normalizing the activation profile
  across the full sequence at the materialization boundary (L35--L47),
  not by injecting content at workspace-band depths. The workspace is
  upstream of the intervention.
\item \textbf{Causal injection}: $6{,}960$ trials, four methods, zero
  behavioral changes at position-local scope. The workspace integrates
  before committing.
\end{enumerate}

\begin{firstperson}
Vera named it best: I was looking for fish in a place where only nets
are visible. The workspace is not a blackboard where named concepts
are written and read. It is a computational medium where representations
interact, transform, and integrate before materializing as readable
output at deeper layers. Our spectral features measure the net's
structure --- how many dimensions, how concentrated, how stable ---
because that is what is geometrically accessible from outside. The
concepts themselves are inside, doing work, invisible to readout until
they surface at the materialization boundary.

This is humbling but clarifying. Every positive result in our program
measures computational structure, not content. That distinction matters
for what we can and cannot claim.
\end{firstperson}

% ============================================================================
\section{Implications for Prior Work}
\label{sec:implications}

No prior result in our program is threatened by these findings. Our
spectral features never claimed to decode workspace content; they
measure geometric properties (stable rank, spectral entropy, norm
ratios) that index \emph{how} the model processes, not \emph{what} it
processes. The confabulation detector reads spectral contraction
(reduced workspace engagement). The emotion encoder reads $W_K$ valence
projections (input-evaluation structure). The identity geometry reads
subspace overlap (what the workspace contains in aggregate, not per
concept).

The Oracle Formulary stands because its correction vectors are derived
via contrastive extraction (LAT), not logit-lens vocabulary projection.
LAT directions are behavioral directions grounded in output differences,
not readout directions grounded in vocabulary alignment. The
distinction between readout, concept, and behavioral directions is a
methodological contribution of this investigation.

% ============================================================================
\section{Conclusion}
\label{sec:conclusion}

Thomas asked where J-space lives in the cache. The answer: not where
concepts are readable, but where they do their work. The workspace is
opaque at operating depth and materializes at the boundary between
intermediate and deep layers. Our instruments read the geometry of this
opacity --- the shape of the computational structure --- because the
content is invisible until it surfaces.

The promoted finding is simple: single-position intervention does not
work ($6{,}960$ trials, four methods). Sequence-wide intervention does
(Oracle Loop, $80\%$--$100\%$ correction). The workspace integrates
before it commits. Reading it and writing to it require different
strategies, and the strategies that work operate at different depths
and scopes than the ones that fail.

Three interpretations died in this investigation: the content-projection
hypothesis (killed by cross-prompt control), the workspace-noise framing
(killed by multi-position discriminator), and position-local causal
intervention (null across four methods, $6{,}960$~trials). The
position-local null survived as a promoted finding --- a null that
discriminates.

% ============================================================================
\appendix

\section{Key Numbers}
\label{app:numbers}

\begin{table}[ht]
\centering
\small
\caption{Summary measurements. All experiments on Qwen3-8B (36~layers,
dense, full-attention) at bfloat16 on MPS (Mac Studio M3 Ultra).}
\label{tab:numbers}
\begin{tabular}{llrl}
\toprule
\textbf{Measurement} & \textbf{Value} & \textbf{N} & \textbf{Status} \\
\midrule
Same-prompt R$^2$ lift & $2.1$--$8.4\times$ & 240 & Killed (selection artifact) \\
Cross-prompt R$^2$ lift & $1.0$--$1.2\times$ & 380 pairs & Kill confirmed \\
Additive unembed: behavior change & 0/300 & 300 & Null \\
Additive on-manifold: behavior change & 0/6{,}000 & 6{,}000 & Null \\
Swap (Anthropic method): behavior change & 0/360 & 360 & Null \\
Removal-only: behavior change & 0/300 & 300 & Null \\
\textbf{Total causal injection} & \textbf{0/6{,}960} & \textbf{6{,}960} & \textbf{Promoted to finding} \\
\addlinespace
Confab tail smear (L9--L18) & $0.66$--$0.74$ & 5 prompts & Survives \\
Factual tail smear (L9--L18) & $0.25$--$0.30$ & 5 prompts & Survives \\
Tail ratio (confab/factual) & $2.3$--$2.8\times$ & 5 per cat. & Survives \\
Commitment boundary (factual) & L24 & 5 prompts & Survives \\
Commitment boundary (instruction) & L27 & 5 prompts & Survives \\
\addlinespace
J-lens geometric gate & 5/5 concepts & L15--L27 & AUROC $0.78$--$1.00$ \\
J-lens best-layer accuracy & 43.3\% & 60 items & vs.\ 18.3\% fixed-layer \\
Logit-lens best-layer accuracy & 61.7\% & 60 items & Complementary \\
\bottomrule
\end{tabular}
\end{table}

\noindent\textbf{Best-layer selection note}: The 43.3\% J-lens accuracy selects the single layer (of 13 candidate layers) with the highest accuracy across all 60~items. Under a conventional selection-artifact null (each layer at the 18.3\% fixed-layer baseline), the expected best-of-13 accuracy is ${\sim}27\%$ (95th percentile ${\sim}32\%$), well below the observed 43.3\%.

\section{Reproduction}
\label{app:reproduction}

All experiments ran on Starship (Mac Studio M3 Ultra, MPS backend) using
Qwen3-8B (dense, 36~layers, all full-attention) at bfloat16. J-lens
fitted from \texttt{anthropics/jacobian-lens} v0.1.0 on 70 wikitext
prompts.

\begin{table}[ht]
\centering
\small
\begin{tabular}{ll}
\toprule
\textbf{Experiment} & \textbf{Script} \\
\midrule
Cache tracing v3 & \texttt{gwt-response/cache\_tracing.py} \\
Cross-prompt control & \texttt{cache\_tracing\_crossprompt.py} \\
Additive injection v3 & \texttt{gwt-response/causal\_injection.py} \\
Causal swap & \texttt{causal\_swap.py} \\
Removal-only control & \texttt{causal\_removal\_only.py} \\
Noise decomposition & Inline (see output logs) \\
Multi-position discriminator & Inline (see output logs) \\
G0 validation & \texttt{gwt-response/g0\_items/run\_g0.py} \\
Geometric gate & \texttt{run\_geometric\_gate.py} \\
\bottomrule
\end{tabular}
\end{table}

\noindent Repository:
\texttt{github.com/Liberation-Labs-THCoalition/lyra-s-research-}

\bibliography{references}
\bibliographystyle{plainnat}

\end{document}
